\chapter{Fundamentação Teórica e Trabalhos Relacionados}\label{cap:fundamentacao_bibliografica}

Este capítulo apresenta os principais conceitos usados no trabalho e alguns estudos relacionados. A discussão aborda percepção visual para drones, visão monocular, estimativa de profundidade, fluxo óptico e navegação reativa. O capítulo também mostra como o problema vem sendo estudado e em que ponto este projeto se diferencia dos trabalhos já publicados.

\section{Navegação autônoma de drones em ambientes complexos}

Quadricópteros conseguem realizar manobras rápidas e operar em locais de difícil acesso. Por isso, são usados em tarefas como inspeção, monitoramento, busca e exploração. Ao mesmo tempo, o sistema de percepção precisa detectar obstáculos e reagir com rapidez, principalmente quando o drone está em movimento.

Em ambientes densos, como florestas, galpões e passagens estreitas, o drone pode encontrar oclusões, objetos finos, mudanças de iluminação, texturas repetitivas e muito movimento aparente na imagem. Esses fatores dificultam a detecção dos obstáculos e a reconstrução do ambiente. Em voos rápidos, não basta detectar corretamente: a informação também precisa chegar a tempo de permitir uma reação \cite{Loquercio2021}.

Abordagens clássicas costumam separar a navegação em etapas de percepção, mapeamento, planejamento e controle. Essa organização facilita o desenvolvimento modular, mas pode introduzir latência e acúmulo de erros, especialmente quando a cena muda rapidamente ou quando o drone se desloca em alta velocidade \cite{Loquercio2021}. Em ambientes desconhecidos, também existe incerteza sobre regiões ainda não observadas, o que pode tornar o planejamento conservador ou arriscado demais \cite{BayesianLearning2015}. Por isso, trabalhos recentes investigam estratégias capazes de aproximar a percepção da tomada de decisão, seja por políticas aprendidas, seja por métricas reativas de risco.

\section{Visão computacional em VANTs}

A visão computacional tem papel central em sistemas autônomos porque permite extrair informações do ambiente a partir de imagens. Em drones, essa percepção visual pode ser usada para detecção de objetos, rastreamento, localização, estimação de movimento, reconstrução do ambiente e prevenção de colisões. No entanto, imagens obtidas por VANTs apresentam desafios específicos, como variação de altitude, ângulo de câmera, movimento rápido, vibração, borramento, oclusões e restrições de processamento embarcado \cite{Ramachandran2021}.

Neste trabalho, detectar um obstáculo não significa classificar cada objeto da cena. O interesse principal é identificar regiões que oferecem risco à trajetória. Assim, é mais importante estimar se algo está próximo ou entrando no caminho do drone do que determinar se o objeto é uma árvore, uma parede ou um poste.

Essa distinção é importante em ambientes naturais, onde os obstáculos podem ser galhos, troncos, folhas ou outros objetos de formas variadas. Nesses casos, o risco depende mais da posição e da proximidade do objeto do que de sua classe. Por esse motivo, o projeto estuda movimento visual e variação de profundidade.

\section{Visão monocular e estimativa de profundidade}

Câmeras monoculares são atrativas para drones por serem leves, baratas e fáceis de integrar. No entanto, estimar profundidade a partir de uma única câmera é um problema difícil, pois a imagem 2D perde diretamente a informação métrica da cena 3D. A profundidade precisa ser inferida por pistas como perspectiva, textura, tamanho relativo, sombreamento, movimento entre quadros e conhecimento aprendido a partir de dados \cite{Vyas2022,Arampatzakis2024}.

Revisões recentes mostram avanços no uso de redes neurais profundas para estimar profundidade monocular \cite{RealTimeMonocular2022,Dhanushkodi2025}. Mesmo assim, o resultado depende do tipo de dado usado no treinamento. Um modelo treinado em ambientes urbanos ou internos pode não funcionar da mesma forma em florestas e parques, onde textura, iluminação e objetos são diferentes.

O custo computacional também deve ser considerado. Em um drone, a estimativa precisa ser produzida dentro do tempo disponível para a tomada de decisão. Dong et al. mostram que existe uma relação entre qualidade da predição, robustez e velocidade de processamento \cite{RealTimeMonocular2022}. Esse ponto é importante para este trabalho porque a reação a um obstáculo precisa ocorrer em intervalos curtos.

A estimativa de profundidade por aprendizado também pode ser tratada como uma distribuição ou medida incerta, em vez de um único valor absoluto por pixel. Liu et al. discutem a importância de modelar incertezas em profundidade monocular, já que diferentes cenas podem produzir ambiguidades visuais semelhantes \cite{Liu2023}. Para o presente projeto, essa observação reforça a escolha de analisar variações temporais e sinais de proximidade, evitando depender apenas de valores absolutos de distância.

\section{Fluxo óptico, movimento visual e proximidade}

O fluxo óptico descreve o deslocamento aparente de pontos ou regiões da imagem entre quadros consecutivos. Em um drone em movimento, esse deslocamento carrega informações sobre a relação entre movimento da câmera, estrutura da cena e proximidade dos objetos. Obstáculos próximos tendem a produzir movimentos aparentes maiores do que regiões distantes, especialmente quando o drone avança em direção a eles.

O método de Lucas--Kanade estima o deslocamento local sob hipóteses de constância de brilho e movimento pequeno \cite{LucasKanade1981}. Pontos com estrutura adequada ao rastreamento podem ser selecionados pelo critério de Shi--Tomasi \cite{ShiTomasi1994}. Como o fluxo observado mistura rotação, translação e geometria da cena, um vetor de grande magnitude não representa, isoladamente, um obstáculo próximo; é necessário compensar o movimento próprio ou combinar o fluxo com outros sinais.

Por isso, o fluxo pode ser útil na navegação reativa. Os vetores podem indicar expansão, aproximação ou movimento lateral antes que seja necessário reconstruir um mapa completo. Trabalhos de odometria visual monocular usam a relação entre o movimento da câmera e a cena para estimar deslocamento e auxiliar a navegação \cite{Tarrio2015,Garcia2016}.

Em aplicações de drones, a frequência e a resolução da imagem também influenciam a robustez da percepção. Godio et al. analisam os efeitos de resolução e frequência em odometria visual-inercial para VANTs, mostrando que o custo computacional e a estabilidade do sistema dependem diretamente da configuração visual utilizada \cite{SimoneGodio2022}. Em voos rápidos, o borramento de movimento também se torna relevante, pois pode degradar a correspondência entre quadros e afetar a qualidade da odometria ou do fluxo visual \cite{Peidong2021}.

Neste trabalho, o fluxo óptico é usado para obter métricas temporais. A análise considera as mudanças entre intervalos consecutivos, em vez de depender apenas de valores de uma imagem isolada. A intenção é representar como o risco muda durante o deslocamento do drone.

Alguns trabalhos usam esse sinal diretamente na prevenção de colisões. Xiao et al. combinam divergência de fluxo e dados inerciais para separar obstáculos reais de mudanças causadas pelo movimento do veículo \cite{Xiao2021}. Hu et al. treinam um modelo de desvio usando o movimento dos pixels de uma câmera \cite{PixelMotion2025}, enquanto Deng et al. usam fluxo translacional e uma máscara de incerteza \cite{Deng2026}. Ng, Li e Kim também analisam a incerteza das correspondências do fluxo \cite{Estimation2109}. Esses trabalhos mostram que o fluxo é útil, mas que deve ser analisado junto a outros sinais, como IMU e movimento do drone.

\section{Desvio reativo e voo em alta velocidade}

A navegação em alta velocidade exige que o drone reduza o tempo entre perceber e agir. Trabalhos como o de Loquercio et al. investigam voo rápido em ambientes complexos por meio de políticas aprendidas que mapeiam observações sensoriais para trajetórias livres de colisão, buscando reduzir a latência dos fluxos tradicionais de processamento \cite{Loquercio2021}. Esse tipo de abordagem mostra que a percepção visual precisa estar fortemente ligada à decisão de movimento quando o objetivo é atravessar ambientes densos com velocidade elevada.

Também existem trabalhos que tratam a navegação segura em ambientes desconhecidos a partir da incerteza sobre colisões futuras. Richter, Vega-Brown e Roy propõem aprendizado bayesiano para navegação rápida e segura, destacando que prever a probabilidade de colisão em regiões ainda não observadas pode permitir velocidades maiores sem abandonar restrições de segurança \cite{BayesianLearning2015}. Embora essa abordagem esteja ligada a planejamento, ela reforça a importância de estimar risco de colisão e não apenas reconstruir o ambiente.

Mais recentemente, sensores visuais alternativos, como câmeras de eventos, também têm sido explorados para desvio de obstáculos em quadricópteros. Bhattacharya et al. mostram que visão monocular baseada em eventos pode apoiar desvio reativo em ambientes internos e externos, aproveitando características como alta faixa dinâmica e menor borramento em movimento rápido \cite{Bhattacharya2024}. Mesmo que este projeto utilize câmera RGB monocular convencional, esse trabalho é relevante por reforçar a relação entre velocidade, percepção visual e reação a obstáculos.

\section{Aprendizado de máquina para percepção visual}

Modelos de aprendizado têm sido usados para estimar profundidade, prever movimento, detectar objetos e aprender políticas de navegação. Em visão monocular, CNNs se destacam por extrair padrões espaciais diretamente da imagem, sendo adequadas para tarefas de profundidade e interpretação de cena \cite{Dhanushkodi2025}. Em navegação autônoma, redes profundas também podem ser usadas para aproximar a relação entre observação visual e ação, como demonstrado em abordagens de voo rápido treinadas em simulação \cite{Loquercio2021}.

Modelos mais simples também podem ser úteis no início do estudo. Uma MLP aplicada a dados tabulares permite verificar se os resumos de fluxo, comandos, movimento e variação visual contêm informação suficiente para estimar mudanças de proximidade. Ela não substitui uma arquitetura convolucional quando é necessário analisar a estrutura completa da imagem, mas torna esta primeira análise mais simples.

Neste projeto, a MLP recebe atributos extraídos e resumidos em valores numéricos. Uma CNN, por outro lado, poderia trabalhar diretamente com as imagens ou com mapas de variação. Essa diferença ajuda a verificar se as métricas agregadas já contêm o sinal necessário ou se parte da informação espacial foi perdida.

Essa escolha não significa que a MLP seja o melhor modelo para os dados. Em bases tabulares de tamanho moderado, modelos baseados em árvores muitas vezes apresentam resultados melhores que redes profundas \cite{OyallonandGaelVaroquaux5220}. A MLP é usada como modelo inicial para testar a representação. O \textit{dropout} e a parada antecipada ajudam a reduzir o sobreajuste \cite{Srivastava2014,Prechelt1998}. A divisão por execução também evita que intervalos muito parecidos da mesma trajetória apareçam no treino e no teste, o que poderia melhorar as métricas de forma artificial \cite{Kapoor2023}.

\section{Simulação e reprodutibilidade}

O Gazebo permite configurar física, cenários e sensores para experimentos robóticos \cite{Koenig7803}, enquanto o PX4 fornece uma arquitetura modular de piloto automático e integração com aplicações externas \cite{Meier8092}. Neste estudo, a combinação é usada em modo \textit{Software in the Loop}: a câmera RGB alimenta a lógica monocular e uma câmera de profundidade alinhada fornece apenas a referência supervisionada.

Mesmo em simulação, diferenças de tempo entre processos, integração numérica e inicialização do estimador podem alterar a trajetória. Como o controlador reage ao estado atual, uma diferença pequena pode aumentar ao longo do voo. Por isso, cada voo completo é tratado como uma execução independente. As execuções de validação e teste são mantidas fixas enquanto o conjunto de treino cresce.

\section{Explicabilidade por atribuição}

Métodos de atribuição baseados em gradiente ajudam a observar quanto a saída ou o erro muda em relação a cada entrada. O método Integrated Gradients calcula essa atribuição ao longo do caminho entre uma referência e a amostra analisada \cite{Sundararajan2017}. Outros trabalhos discutem suas propriedades e o problema de regiões em que o gradiente fica saturado \cite{Walker5289,Axiomatic2025}.

Como o modelo deste trabalho recebe atributos tabulares, não se aplica Grad-CAM, técnica voltada a ativações espaciais de redes convolucionais. A análise retropropaga o erro até as entradas padronizadas e agrega o módulo dos gradientes. O resultado representa sensibilidade local, não causalidade. Quando o gradiente é exatamente zero, não há ranking local válido e a situação deve ser explicitada em vez de associar nomes arbitrários a importâncias nulas.

\section{Posicionamento deste trabalho}

Os estudos revisados mostram que a navegação em ambientes densos depende da percepção visual, do tempo de resposta e do tratamento da incerteza. Trabalhos sobre voo rápido procuram reduzir a demora entre percepção e ação \cite{Loquercio2021,Bhattacharya2024}. Estudos de profundidade monocular mostram tanto as possibilidades quanto as limitações de usar apenas uma câmera \cite{Vyas2022,RealTimeMonocular2022}. Já os trabalhos de odometria e fluxo indicam que o movimento entre quadros pode fornecer informações úteis para a navegação \cite{Tarrio2015,SimoneGodio2022}.

O objetivo deste projeto não é apresentar o primeiro método de desvio monocular por fluxo óptico, pois esse tipo de abordagem já aparece em outros trabalhos \cite{Xiao2021,PixelMotion2025}. A proposta é montar um processo de análise que reúna fluxo esparso, IMU, movimento, comandos e profundidade de referência. Os dados são representados como variações entre intervalos, e o aumento do conjunto de treino é avaliado mantendo validação e teste fixos. Também são analisadas as entradas que mais influenciam os erros do modelo. A profundidade simulada é usada apenas como referência; a navegação continua baseada nos sinais monoculares.
